% Source: https://artofproblemsolving.com/community/c3h2063723_xy__x__y__232
% Stable identity: aops-topic-2063723-q01
Adding $1$: $(x+1)(y+1)=2^{32}+1$. Now
\[
2^{32}+1=641\cdot6700417 .
\]
(Indeed $641=5\cdot2^7+1=5^4+2^4$. From $5\cdot2^7\equiv-1\pmod{641}$ we get $5^4\cdot2^{28}\equiv1$, and $5^4\equiv-2^4$, so $-2^{32}\equiv1$, i.e.\ $641\mid2^{32}+1$; the quotient is $6700417$.) Both factors are prime ($641$ by trial division up to $25$; $6700417$ as given), so the divisors of $2^{32}+1$ are $1$, $641$, $6700417$ and $2^{32}+1$.

Since $x,y\ge1$, both $x+1$ and $y+1$ are at least $2$, so $\{x+1,y+1\}=\{641,6700417\}$. The solutions are
\[
\boxed{(x,y)=(640,\,6700416)\ \text{ and }\ (6700416,\,640)}.
\]
(If $0$ were allowed as a natural number, $(0,2^{32})$ and $(2^{32},0)$ would also work.)
